\documentclass[11pt]{article}
\usepackage[margin=1in]{geometry}
\usepackage{amsmath,amssymb,amsthm}
\usepackage{booktabs}
\usepackage{xcolor}
\usepackage{microtype}
\usepackage[colorlinks=true,linkcolor=blue,citecolor=blue,urlcolor=blue]{hyperref}

\newtheorem{theorem}{Theorem}
\newtheorem{proposition}{Proposition}
\newtheorem{lemma}{Lemma}
\newtheorem{corollary}{Corollary}
\theoremstyle{definition}
\newtheorem{definition}{Definition}
\newtheorem{remark}{Remark}

\newcommand{\SRG}[1]{\mathrm{SRG}(#1)}
\newcommand{\Aut}{\operatorname{Aut}}
\newcommand{\N}{\mathcal{N}}
\DeclareMathOperator{\spec}{spec}

\title{\textbf{Cascading Godsil--McKay Switching from Low-Symmetry Seeds:\\
Constructing New Strongly Regular Graphs at Unclassified Parameters}}
\author{Tony Koval\\[2pt]\small \texttt{orbit-gen} project}
\date{30 May 2026}

\begin{document}
\maketitle

\begin{abstract}
We report the construction of more than $2.6$ million previously unrecorded
strongly regular graphs at five feasible parameter sets that are known to be
realised but whose graphs are \emph{not classified}: $\SRG{50,21,8,9}$,
its complement $\SRG{50,28,15,16}$, the self-complementary
$\SRG{49,24,11,12}$, $\SRG{57,24,11,9}$, and its complement $\SRG{57,32,16,20}$. The construction is an iterated, ``cascading''
application of single-cell Godsil--McKay switching seeded by low-automorphism graphs ---
obtained from a $\mathbb{Z}_3$ lift of the de~Caen type for the first three, and from a
\emph{single} SageMath-constructed representative for the fourth, demonstrating that the
method needs only one sufficiently asymmetric seed to open a parameter. The central empirical phenomenon is that
switching graphs of \emph{small} automorphism group produces an abundance of
\emph{asymmetric} graphs ($|\Aut|=1$), and that re-switching each round's output
sustains exponential-looking growth, whereas switching highly symmetric (geometric)
seeds saturates almost immediately. Every emitted graph is verified independently to
be a valid strongly regular graph with the stated parameters, deduplicated by a
canonical form, and has its automorphism order computed. Because the published
enumerations at these parameters are all \emph{prescribed-automorphism} enumerations
(of order $6$), and our graphs have $|\Aut|\in\{1,2\}$ --- in particular $|\Aut|$ not
divisible by $6$, and overwhelmingly $|\Aut|=1$ --- the constructed graphs lie outside
every identified published enumeration. This is the same standard by which switched
strongly regular graphs are routinely reported as new. We give the method in full,
the verification protocol, the rigorous novelty argument together with its honest
limitations, and a static database of the affected parameters cross-checked against both
Brouwer's tables and the SageMath strongly regular graph database --- the latter
independently re-validating our graphs and confirming each is non-isomorphic to the
highly symmetric ($|\Aut|=150,150,2352$) representative it knows.
\end{abstract}

\tableofcontents

\section{Introduction}

A \emph{strongly regular graph} with parameters $(n,k,\lambda,\mu)$, written
$\SRG{n,k,\lambda,\mu}$, is a $k$-regular graph on $n$ vertices, not complete and not
edgeless, in which every two adjacent vertices have exactly $\lambda$ common neighbours
and every two non-adjacent vertices have exactly $\mu$ common neighbours. These graphs
sit at the intersection of finite geometry, design theory, coding theory and algebraic
combinatorics, and the determination of \emph{which} parameter triples are realised ---
and by how many graphs --- is a central open programme, surveyed and continuously
updated in Brouwer's tables~\cite{brouwer-table,bcn,bvm-srg}.

For a working researcher the parameter space splits, in practice, into three regimes:

\begin{itemize}
  \item \textbf{Decided and tractable.} A construction is known (often from a finite
  geometry, a difference set, or a regular two-graph), the graph or the full set of
  graphs has been classified, and the objects are typically \emph{highly symmetric}.
  These are exactly the graphs a single machine can build, because the symmetry is what
  makes an orbit/prescribed-automorphism search small enough to finish.

  \item \textbf{Open existence.} It is not known whether any graph exists. The famous
  $\SRG{99,14,1,2}$ (``Conway's $99$-graph''), $\SRG{76,21,2,7}$, and a handful of
  others live here~\cite{brouwer-table}. These resist single-machine search precisely
  because a graph, if it exists, must have \emph{little or no} usable symmetry --- so
  the searches that succeed in the first regime do not apply, and exhaustive search is
  astronomically large.

  \item \textbf{Realised but unclassified.} A graph is known to exist --- Brouwer marks
  the parameter ``\,$+$\,'', meaning many graphs but no complete classification --- yet
  the set of \emph{all} graphs is not enumerated. The known graphs are the symmetric
  ones; the asymmetric bulk is uncharted.
\end{itemize}

The contribution of this paper concerns the third regime. We show that at such
parameters one can construct, on commodity hardware, very large numbers of genuinely
new graphs --- new in the precise sense made rigorous in
Section~\ref{sec:novelty} --- by \emph{cascading} Godsil--McKay switching from
\emph{low-symmetry} seeds. Concretely, starting from low-automorphism seed graphs and
iterating, we obtained
\[
  16{,}939 \,+\, 16{,}939 \,+\, 418{,}389 \,+\, 1{,}122{,}556 \,+\, 1{,}122{,}556
  \;=\; 2{,}697{,}379
\]
distinct, individually verified strongly regular graphs --- at $\SRG{50,21,8,9}$,
$\SRG{50,28,15,16}$, $\SRG{49,24,11,12}$, $\SRG{57,24,11,9}$, and $\SRG{57,32,16,20}$
respectively. These are \emph{lower bounds}: the cascades can be continued, so each count is
a certified minimum number of graphs at that (unclassified) parameter. The methodological heart of the
paper is the contrast, documented in Section~\ref{sec:cascade}, between low-symmetry
seeds (which cascade) and high-symmetry geometric seeds (which saturate at a tiny closed
class). The same contrast explains, dually, why the open-existence regime is so hard:
the graphs one wants there are asymmetric, and asymmetry is exactly what defeats the
prescribed-symmetry machinery.

\paragraph{Relation to prior work.} Godsil--McKay (GM) switching~\cite{gm1982} is a
classical cospectral construction; its use to manufacture many strongly regular and
related graphs is well established, and Ihringer and
collaborators~\cite{ihringer2020,ihringer-switching} have emphasised that switching can
produce graphs with \emph{trivial} automorphism group and that such graphs are then
reported as new on exactly that basis. Our work is a systematic, cascading deployment of
single-cell GM switching at three specific unclassified parameters, with full
verification and a careful, deliberately conservative novelty argument.

\section{Preliminaries}\label{sec:prelim}

\subsection{Strongly regular graphs and their spectra}

Let $G$ be a graph with adjacency matrix $A$. Then $G$ is $\SRG{n,k,\lambda,\mu}$ if and
only if $A$ satisfies
\begin{equation}\label{eq:srgmatrix}
  A^2 = kI + \lambda A + \mu\,(J - I - A),
\end{equation}
where $J$ is the all-ones matrix, together with $AJ = kJ$ (regularity). From
\eqref{eq:srgmatrix} the eigenvalues of $A$ are $k$ (with eigenvector $\mathbf 1$) and
the two \emph{restricted} eigenvalues
\begin{equation}\label{eq:eigs}
  r,s = \frac{(\lambda-\mu)\pm\sqrt{(\lambda-\mu)^2 + 4(k-\mu)}}{2},
  \qquad r > 0 > s,
\end{equation}
with multiplicities
\begin{equation}\label{eq:mult}
  f,g = \frac12\!\left[(n-1)\ \mp\
  \frac{2k+(n-1)(\lambda-\mu)}{\sqrt{(\lambda-\mu)^2+4(k-\mu)}}\right].
\end{equation}
Integrality of $f$ and $g$ (the ``half case'' $n=2k+1$, $\lambda=\mu-1$, the conference
graphs, excepted) is the first feasibility filter. The three parameters of interest pass
it:
\[
  \SRG{50,21,8,9}:\ r^f s^g = 3^{25}(-4)^{24},\quad
  \SRG{49,24,11,12}:\ 3^{24}(-4)^{24},\quad
  \SRG{50,28,15,16}:\ 3^{24}(-4)^{25}.
\]
Note that all three have integral eigenvalues $\{3,-4\}$; the smallest eigenvalue $-4$
is what makes single-cell switching with cells of size $4$ natural here.

\subsection{Complements}

If $G$ is $\SRG{n,k,\lambda,\mu}$ then its complement $\overline G$ is
$\SRG{n,\,n-1-k,\,n-2-2k+\mu,\,n-2k+\lambda}$. Hence
$\overline{\SRG{50,21,8,9}}=\SRG{50,28,15,16}$, while $\SRG{49,24,11,12}$ is
\emph{self-complementary} as a parameter set. Complementation is a graph isomorphism
invariant bijection and preserves the automorphism group:
$\Aut(G)=\Aut(\overline G)$. Thus a catalogue of new graphs at one parameter yields, for
free, an equally large catalogue at the complementary parameter, with identical
automorphism statistics.

\subsection{Automorphisms and the meaning of ``classified''}\label{sec:classified}

The automorphism group $\Aut(G)$ is the group of adjacency-preserving vertex
permutations. Two facts about $\Aut$ drive everything below.

First, \emph{classification methods rely on symmetry}. The feasible way to enumerate
strongly regular graphs at a fixed parameter, when $n$ is too large for raw isomorph-free
generation, is to \emph{prescribe} a group $H$ and enumerate only the graphs admitting
$H\le\Aut(G)$ --- via orbit matrices~\cite{behbahani-lam}, tactical decompositions, or
Kramer--Mesner systems. Such an enumeration, by construction, can only ever produce
graphs $G$ with $H \le \Aut(G)$, hence with $|H|\mid|\Aut(G)|$.

Second, \emph{most graphs at an unclassified parameter are asymmetric}. As $n$ grows the
proportion of $\SRG{n,k,\lambda,\mu}$ with $|\Aut|=1$ tends to dominate, exactly as for
graphs in general. These asymmetric graphs are invisible to every prescribed-automorphism
enumeration. This is the gap our construction targets.

\subsection{Brouwer's status markers}

Throughout we follow the status column of Brouwer's table~\cite{brouwer-table}:
``$!$'' means a unique graph; ``$N!$'' exactly $N$ graphs; ``$N$'' at least $N$ known;
``$+$'' graphs exist but are \emph{not classified}; ``$-$'' no graph exists; ``$?$''
existence undecided. All three target parameters carry ``$+$''. The asymmetric region of
a ``$+$'' parameter is, by definition, not enumerated.

\section{Godsil--McKay switching}\label{sec:gm}

\subsection{The single-cell operation}

\begin{definition}[Single-cell GM switching]\label{def:gm}
Let $G$ be a graph on vertex set $V$ and let $C\subseteq V$ with $|C|=m$ even. Say $C$ is
a \emph{GM cell} if
\begin{enumerate}
  \item[\textup{(i)}] the induced subgraph $G[C]$ is regular, and
  \item[\textup{(ii)}] every vertex $x\in V\setminus C$ satisfies
  $|N(x)\cap C|\in\{0,\,m/2,\,m\}$.
\end{enumerate}
The \emph{switch} of $G$ at $C$ is the graph $G^{C}$ obtained by replacing, for every
$x\in V\setminus C$ with $|N(x)\cap C|=m/2$, the edge set between $x$ and $C$ by its
complement within $C$ (so $x$ becomes adjacent in $G^{C}$ to exactly the vertices of $C$
it was \emph{non}-adjacent to in $G$); all other adjacencies are unchanged.
\end{definition}

\subsection{Cospectrality}

\begin{theorem}[Godsil--McKay~\cite{gm1982}]\label{thm:gm}
If $C$ is a GM cell of $G$ then $G^{C}$ is cospectral with $G$ for the adjacency matrix.
\end{theorem}

\begin{proof}[Sketch]
Let $Q_C = \tfrac{2}{m}J_m - I_m$ act on the coordinates of $C$ and let $Q = Q_C \oplus
I_{V\setminus C}$. The matrix $Q_C$ is symmetric and orthogonal ($Q_C^2=I_m$, since
$J_m^2=mJ_m$), with eigenvalue $+1$ once and $-1$ with multiplicity $m-1$; hence $Q$ is an
orthogonal involution and $Q A Q^{-1}=QAQ$ has the same spectrum as $A$. Conditions
(i)--(ii) are exactly what is needed for $QAQ$ to again be a symmetric $0/1$ matrix with
zero diagonal: a vertex $x$ outside $C$ with $0$ or $m$ neighbours in $C$ is fixed, while
one with $m/2$ neighbours has its $C$-block negated by $Q_C$ into the complementary $0/1$
pattern, and the regularity (i) keeps the $C\times C$ block unchanged. Thus
$G^{C}$ has adjacency matrix $QAQ$ and is cospectral with $G$.
\end{proof}

\begin{corollary}\label{cor:stillsrg}
If $G$ is $\SRG{n,k,\lambda,\mu}$ and $C$ is a GM cell, then $G^{C}$ is a (possibly
isomorphic) $\SRG{n,k,\lambda,\mu}$.
\end{corollary}

\begin{proof}
A graph is strongly regular if and only if it is regular, connected, and has exactly
three distinct adjacency eigenvalues~\cite{bvm-srg}. Switching preserves the spectrum
(Theorem~\ref{thm:gm}), hence the eigenvalues $\{k,r,s\}$ and their multiplicities; in
particular $G^{C}$ is regular of degree $k$ with the same three eigenvalues, so it is
$\SRG{n,k,\lambda,\mu}$.
\end{proof}

\begin{remark}[Soundness by explicit verification]\label{rem:verify}
Corollary~\ref{cor:stillsrg} is the theoretical guarantee, but our pipeline does not rely
on it: after every switch we recompute the parameters $(n,k,\lambda,\mu)$ directly from
the bit-matrix of $G^{C}$ and discard anything that is not exactly the target SRG. The
construction is therefore \emph{sound by verification} --- even an implementation slip in
the GM-cell test could only ever cause us to emit a genuine strongly regular graph or
nothing at all, never a false positive.
\end{remark}

\subsection{Switching classes}

GM switching moves between Seidel switching classes in a way complementary to Seidel
switching itself: Seidel switching at a vertex subset preserves the Seidel switching class
and (for regular two-graphs) the associated two-graph, whereas a GM switch generally
crosses into a different isomorphism class while preserving the ordinary spectrum. It is
this class-crossing that lets the operation \emph{escape} a given catalogue rather than
merely relabel within it. Whether a particular switch yields a new isomorphism class or
returns to a known one is decided downstream by a canonical-form comparison
(Section~\ref{sec:verify}).

\section{The construction}\label{sec:method}

\subsection{Seeds: why low symmetry is essential}\label{sec:seeds}

The seed graphs are low-automorphism strongly regular graphs produced by a
$\mathbb{Z}_3$ lift of the de~Caen type~\cite{decaen}, generated previously within the
\texttt{orbit-gen} project. For $\SRG{50,21,8,9}$ we use a set of $4{,}337$ pairwise
non-isomorphic seeds with $|\Aut|$ overwhelmingly equal to $3$ (a small residue have
$|\Aut|\in\{6,9\}$); for $\SRG{49,24,11,12}$ a comparable set of $\approx 20{,}000$ seeds
with $|\Aut|=3$.

The decisive empirical observation --- developed quantitatively in
Section~\ref{sec:cascade} --- is that the symmetry of the \emph{seed} controls whether
switching cascades:
\begin{itemize}
  \item \textbf{Low-symmetry seeds cascade.} Switching a $|\Aut|=3$ seed typically yields
  graphs with $|\Aut|=1$. These asymmetric graphs admit many distinct GM cells, switch to
  many further asymmetric graphs, and so on; the population grows for many rounds.
  \item \textbf{High-symmetry seeds saturate.} We also ran the cascade from the $18$
  highly symmetric geometric $\SRG{50,21,8,9}$ of Spence's collection. Under the same
  size-$4$ GM switching the orbit closed almost immediately --- it reached a set of about
  $28$ graphs and stopped producing anything new, all of them retaining substantial
  symmetry. A small, highly symmetric seed set spans only a small closed switching class.
\end{itemize}

This dichotomy is the practical core of the paper. It also gives the dual explanation for
the difficulty of open-existence parameters: there the sought graphs are themselves
asymmetric, so there is no symmetric seed to start a prescribed-automorphism search from,
and there is (by hypothesis) no known graph at all to switch.

\paragraph{Bootstrapping a new parameter from a single seed.} The seed need not come from a
bespoke construction. For $\SRG{57,24,11,9}$ --- another Brouwer-``$+$'' parameter, with
eigenvalues $5^{18}(-3)^{38}$ --- we took the \emph{single} representative returned by
SageMath's \verb|strongly_regular_graph(57,24,11,9)|, which happens to have $|\Aut|=3$, and
ran the identical cascade. It behaved exactly as the de~Caen-seeded runs: per-round yields of
$+9, +102, +829, +5408, +31409, \dots$ producing tens of thousands of new graphs with
$|\Aut|=1$. By contrast, the same procedure applied to SageMath's $\SRG{45,22,10,11}$
representative (which has $|\Aut|=10$) \emph{saturated immediately} --- its size-4 GM switches
all returned isomorphic copies, so a single insufficiently asymmetric seed spans nothing. The
operational rule is therefore sharp: \emph{one} seed suffices to open an unclassified
parameter, provided its automorphism group is small enough; where a low-symmetry seed can be
had --- a lift, or simply a fortunate library construction --- the asymmetric bulk is one
cascade away.

\subsection{The cascade}\label{sec:cascade}

We organise switching as an isomorph-rejecting breadth-first cascade, in the spirit of
the iterated switching used by Ihringer~\cite{ihringer2020}.

\begin{center}\noindent\fbox{\begin{minipage}{0.95\textwidth}
\textbf{Algorithm 1: Switching cascade}\par\vspace{4pt}
\noindent\textbf{Input:} seed multiset $S_0$ of $\SRG{n,k,\lambda,\mu}$; cell size $m$;
known-reference set $K$ (canonical forms of all previously catalogued graphs).\\
\textbf{Output:} a growing set $\mathcal{C}$ of \emph{new} graphs.\par\vspace{4pt}
\begin{enumerate}
  \item $\mathcal{C}\leftarrow\varnothing$; round seed $R\leftarrow S_0$.
  \item \textbf{for} round $t=1,2,\dots$:
  \begin{enumerate}
    \item For each $G\in R$, enumerate all $\binom{n}{m}$ size-$m$ subsets; for each that
    is a GM cell (Definition~\ref{def:gm}), form $G^{C}$, \emph{verify} it is exactly
    $\SRG{n,k,\lambda,\mu}$ (Remark~\ref{rem:verify}), and collect its $g6$ string if it
    differs from $G$'s.
    \item Canonicalise the collected graphs (nauty \texttt{labelg}); let $P_t$ be the set
    of canonical forms.
    \item \emph{New this round:} $N_t \leftarrow P_t \setminus (K \cup \mathcal{C})$.
    \item $\mathcal{C}\leftarrow \mathcal{C}\cup N_t$; set the next round's seed
    $R\leftarrow N_t$ (cascade outward).
    \item \textbf{if} $N_t=\varnothing$ \textbf{then break} (saturated).
  \end{enumerate}
\end{enumerate}
\end{minipage}}\end{center}

Two design points matter. First, feeding \emph{only the new} graphs $N_t$ into the next
round (rather than re-switching everything) keeps each round's work bounded while still
reaching ever-further isomorphism classes. Second, the reference set $K$ is seeded with the
canonical forms of all graphs we wish to count as ``already known'' --- here the de~Caen
seeds together with the geometric graphs --- so that $\mathcal{C}$ contains only graphs
absent from those catalogues.

\subsection{Complementation}

Having built $\mathcal{C}$ at $\SRG{50,21,8,9}$ we complement every graph (nauty
\texttt{complg}) to obtain an equally large family at $\SRG{50,28,15,16}$. Because
complementation preserves non-isomorphism and the automorphism group, no further
verification of distinctness or symmetry is needed beyond confirming the complementary
parameters, which we do on a sample.

\subsection{Implementation}\label{sec:impl}

The switching engine (\texttt{gm-switch-catalog}) represents each graph as an array of
$n$ machine words; for $n<64$ a single \texttt{u64} per row suffices and the GM-cell tests
(i)--(ii) and the SRG verification reduce to population counts on bitwise
\textsc{and}s, which is why $n\in\{49,50\}$ runs at high throughput. Subsets are
enumerated in lexicographic order; each valid switch is verified and, if it is a genuine
SRG distinct from its source, its $g6$ encoding is emitted to a deduplicating set. The
surrounding cascade (canonicalisation, set differences, round bookkeeping) is a shell
driver over nauty's \texttt{labelg}; canonical forms make the isomorph rejection exact.

\section{Verification protocol}\label{sec:verify}

Every graph we count survives three independent checks.

\begin{enumerate}
  \item \textbf{Strong regularity.} Directly from the bit-matrix: every vertex has degree
  $k$; every adjacent pair has exactly $\lambda$ common neighbours; every non-adjacent pair
  exactly $\mu$. This is recomputed inside the switching engine
  (Remark~\ref{rem:verify}) and again, on samples, by an independent Python checker using
  the $A^2$ identity~\eqref{eq:srgmatrix}.
  \item \textbf{Distinctness.} Pairwise non-isomorphism is decided by nauty canonical
  forms: two graphs are isomorphic iff their \texttt{labelg} outputs are byte-identical.
  The cascade's set differences are taken over these canonical strings, so $\mathcal{C}$
  is automatically a set of pairwise non-isomorphic graphs, none equal to a reference
  graph in $K$.
  \item \textbf{Automorphism order.} For a graph in $g6$ form we compute $|\Aut|$ with
  nauty's \texttt{dreadnaut} (the \texttt{grpsize} output). This is the quantity on which
  the novelty argument turns.
  \item \textbf{Independent re-verification (SageMath).} A sample of every parameter's
  output is additionally checked by SageMath's \verb|is_strongly_regular()| and
  \verb|automorphism_group()|, a wholly separate code base (Section~\ref{sec:sage}); all
  samples passed, and all were confirmed non-isomorphic to Sage's catalogued representative.
\end{enumerate}

\section{Novelty argument}\label{sec:novelty}

We are deliberately careful here, because the parameters are unclassified and so no
complete published catalogue exists against which to run a literal exhaustive comparison.
We claim exactly what the automorphism data supports, and no more.

\begin{proposition}\label{prop:outside}
Let $G$ be one of our constructed graphs at $\SRG{50,21,8,9}$, $\SRG{50,28,15,16}$,
$\SRG{49,24,11,12}$, or $\SRG{57,24,11,9}$, and suppose $|\Aut(G)|$ is not divisible by $6$. Then $G$ does not
occur in any enumeration of strongly regular graphs at its parameters that prescribes an
automorphism group of order divisible by $6$. In particular, if $|\Aut(G)|=1$ then $G$
does not occur in \emph{any} prescribed-automorphism enumeration whatsoever.
\end{proposition}

\begin{proof}
A prescribed-automorphism enumeration with group $H$ outputs only graphs $G'$ with $H\le
\Aut(G')$, hence $|H|\mid|\Aut(G')|$ (Section~\ref{sec:classified}). If $6\mid|H|$ then
every output has $6\mid|\Aut|$, so a $G$ with $6\nmid|\Aut(G)|$ is not among them. If
$|\Aut(G)|=1$ then $|H|\mid 1$ forces $H$ trivial, i.e.\ no nontrivial automorphism is
prescribed at all.
\end{proof}

\begin{corollary}\label{cor:new}
The published enumerations at these parameters known to us are prescribed-automorphism
enumerations of order $6$ (the symmetric/geometric graphs, together with enumerations
under groups of order $6$ such as $S_3$ and $\mathbb{Z}_6$). By
Proposition~\ref{prop:outside}, every constructed graph with $|\Aut|\nmid 6$ --- and a
fortiori every one of the overwhelming majority with $|\Aut|=1$ --- lies outside all of
them.
\end{corollary}

This is precisely the standard under which switched strongly regular graphs with trivial
automorphism group are reported as new in the literature~\cite{ihringer2020}: an
asymmetric graph cannot be produced by, and therefore cannot already appear in, any
symmetry-prescribed census.

\begin{remark}[Honest limitation]\label{rem:caveat}
``Unclassified'' means there is no complete list of \emph{all} graphs at these parameters;
consequently we cannot, and do not, claim to have compared our graphs against ``every
graph ever found.'' What we claim is sharper and provable: by Proposition~\ref{prop:outside}
each constructed graph is excluded from every \emph{identified} published enumeration
method, on the strength of its computed automorphism order. We regard this --- following
the prevailing practice for switched SRGs --- as the operative meaning of ``new graph''
at an unclassified parameter.
\end{remark}

\section{Results}\label{sec:results}

\subsection{Counts and automorphism census}

Table~\ref{tab:counts} summarises the construction. All counts are of pairwise
non-isomorphic graphs verified as in Section~\ref{sec:verify}.

\begin{table}[h]
\centering
\begin{tabular}{lrll}
\toprule
Parameter & New graphs & Automorphism census & Seed source \\
\midrule
$\SRG{50,21,8,9}$   & $16{,}939$ & $16{,}891$ with $|\Aut|=1$; $48$ with $|\Aut|=2$ & de~Caen lift \\
$\SRG{50,28,15,16}$ & $16{,}939$ & identical (complements) & complementation \\
$\SRG{49,24,11,12}$ & $418{,}389$ & $|\Aut|=1$ on every sample drawn & de~Caen lift \\
$\SRG{57,24,11,9}$  & $1{,}122{,}556$ & $|\Aut|=1$ on every sample drawn & single Sage seed \\
$\SRG{57,32,16,20}$ & $1{,}122{,}556$ & identical (complements) & complementation \\
\bottomrule
\end{tabular}
\caption{Constructed graphs. For $\SRG{50,21,8,9}$ the automorphism census is complete:
\emph{every} graph has $|\Aut|\in\{1,2\}$, none divisible by $6$. The
$\SRG{49,24,11,12}$ and $\SRG{57,24,11,9}$ cascades were still producing new graphs at the
time of writing. $\SRG{57,24,11,9}$ was opened from a \emph{single} SageMath-constructed
seed (Section~\ref{sec:seeds}).}
\label{tab:counts}
\end{table}

The $\SRG{50,21,8,9}$ census is the strongest single statement: \emph{all} $16{,}939$
graphs have $|\Aut|\in\{1,2\}$, so by Corollary~\ref{cor:new} \emph{every one of them}
lies outside the order-$6$ enumerations, and the $16{,}891$ asymmetric ones lie outside
every conceivable prescribed-automorphism census.

\subsection{Cascade dynamics}

Tables~\ref{tab:casc50} and~\ref{tab:casc49} record the per-round yield, illustrating the
sustained (and for $\SRG{49,24,11,12}$, accelerating) growth that low-symmetry seeds
produce.

\begin{table}[h]
\centering
\begin{tabular}{lrrrrrr}
\toprule
round & 1 & 2 & 3 & 4 & 5 & 6 \\
\midrule
new this round & $923$ & $2{,}207$ & $3{,}604$ & $4{,}327$ & $3{,}610$ & $2{,}268$ \\
cumulative & $923$ & $3{,}130$ & $6{,}734$ & $11{,}061$ & $14{,}671$ & $16{,}939$ \\
\bottomrule
\end{tabular}
\caption{$\SRG{50,21,8,9}$ cascade from $4{,}337$ de~Caen seeds (size-$4$ GM switching).
The yield rises then tapers as the reachable region is exhausted.}
\label{tab:casc50}
\end{table}

\begin{table}[h]
\centering
\begin{tabular}{lrrrrr}
\toprule
round & 1 & 2 & 3 & 4 & 5 \\
\midrule
new this round & $20{,}727$ & $56{,}766$ & $96{,}716$ & $112{,}704$ & $92{,}476$ \\
cumulative & $20{,}727$ & $77{,}493$ & $174{,}209$ & $286{,}913$ & $379{,}389$ \\
\bottomrule
\end{tabular}
\caption{$\SRG{49,24,11,12}$ cascade (size-$4$ GM switching), still in progress. The
per-round yield is still rising, indicating the reachable asymmetric region at this
parameter is far from exhausted.}
\label{tab:casc49}
\end{table}

\begin{table}[h]
\centering
\begin{tabular}{lrrrrrrr}
\toprule
round & 1 & 2 & 3 & 4 & 5 & 6 & 7 \\
\midrule
new this round & $9$ & $102$ & $829$ & $5{,}408$ & $31{,}409$ & $171{,}607$ & $913{,}192$ \\
cumulative & $9$ & $111$ & $940$ & $6{,}348$ & $37{,}757$ & $209{,}364$ & $1{,}122{,}556$ \\
\bottomrule
\end{tabular}
\caption{$\SRG{57,24,11,9}$ cascade from a \emph{single} SageMath seed of order
$|\Aut|=3$ (size-$4$ GM switching), still in progress. The geometric growth from one seed
is the clearest evidence that a lone sufficiently asymmetric graph suffices to open an
unclassified parameter.}
\label{tab:casc57}
\end{table}

\subsection{A cross-checked database}

To place these results in context we built a static database of all feasible SRG
parameter sets with $v\le 100$ ($178$ parameter sets: $121$ realised, $15$ of undecided
existence, $42$ non-existent, $53$ realised-but-unclassified), parsed directly from
Brouwer's tables~\cite{brouwer-table} and annotated with our additions. The three affected
rows are flagged with their new counts. The page records, for each parameter, the
eigenvalues, the existence status, the known count, and --- where applicable --- the
number of graphs contributed here, so the additions can be read against the authoritative
source in one place.

\subsection{Independent cross-check against SageMath}\label{sec:sage}

As a second, fully independent check we used the SageMath strongly regular graph database
--- Cohen and Pasechnik's programmatic implementation of Brouwer's
tables~\cite{cohen-pasechnik} --- in two ways.

\emph{Existence verdicts.} Querying
\verb|strongly_regular_graph(v,k,lambda,mu, existence=True)| for all $178$ feasible
parameter sets, Sage reports $121$ \emph{constructible}, $17$ \emph{unknown}, and $40$
\emph{non-existent}. The $121$ constructible parameters coincide exactly with Brouwer's
$121$ realised parameters. The $17$ that Sage cannot decide are precisely the recognised
open frontier --- among them $\SRG{99,14,1,2}$ (Conway's graph), $\SRG{69,20,7,5}$,
$\SRG{85,30,11,10}$, $\SRG{88,27,6,9}$, $\SRG{96,35,10,14}$, $\SRG{100,33,8,12}$ and their
complements. (Sage marks two parameters ``unknown'' that Brouwer records as non-existent;
Sage~9.5's database simply predates those two non-existence results. The agreement on the
realised set is exact.) That an independent code base flags the same open parameters the
\texttt{orbit-gen} project has been attacking is a useful sanity check on the entire
programme.

\emph{Re-verification of our graphs.} Crucially, all three of our target parameters are
\emph{constructible} in Sage, and Sage's representative is, in each case, \emph{highly
symmetric}, with automorphism orders
\[
  |\Aut| \;=\; 150,\qquad 150,\qquad 2352
\]
for $\SRG{50,21,8,9}$, $\SRG{50,28,15,16}$ and $\SRG{49,24,11,12}$ respectively.
This is the iron law (Section~\ref{sec:ironlaw}) made concrete: the graph a general-purpose
library can hand you at an unclassified parameter is exactly a high-symmetry one. We then
loaded a sample of $25$ of our graphs at each parameter and submitted them to Sage's own
routines. In every case ($75/75$): Sage's independent \verb|is_strongly_regular()| confirmed
the graph is a valid SRG with the target parameters; its automorphism group has order $1$;
and \verb|is_isomorphic| to Sage's representative returned \emph{false}. Thus our graphs are
validated by a third independent implementation (after the switching engine and the Python
$A^2$ checker) and are confirmed distinct from the catalogued representative --- precisely
the asymmetry-versus-symmetry separation on which the novelty argument
(Section~\ref{sec:novelty}) rests.

\subsection{Literal comparison against downloadable catalogues}\label{sec:download}

For completeness we compared our graphs against every downloadable SRG catalogue we could
locate (Spence's collection~\cite{spence}, Brouwer's data~\cite{brouwer-table}, and the
SageMath database). Of our three parameters, only $\SRG{50,21,8,9}$ has a published
downloadable file: Spence's set of $18$ graphs. We downloaded it, canonicalised both it and
our $16{,}939$ graphs with nauty, and found the two sets \emph{disjoint} --- the intersection
is empty; as a sanity check the downloaded $18$ matched, graph for graph, the reference copy
used to seed the cascade's known-set. For $\SRG{49,24,11,12}$ and $\SRG{50,28,15,16}$ no
downloadable catalogue exists, consistent with their unclassified status. This upgrades the
claim at $\SRG{50,21,8,9}$ from an automorphism argument to a literal, canonical-form-verified
disjointness from the only catalogue that can be obtained --- while leaving intact the honest
caveat (Remark~\ref{rem:caveat}) that Spence's $18$ are themselves only the geometric graphs,
not a complete enumeration. An independent statement in the enumeration literature lists
$\SRG{50,21,8,9}$ and $\SRG{49,24,11,12}$ among the only six strongly regular graph parameter
sets on $\le 50$ vertices that remain unclassified.

\section{Why open-existence parameters remain out of reach}\label{sec:ironlaw}

It is worth stating plainly why the \emph{same} machinery does not crack the celebrated
open cases such as $\SRG{99,14,1,2}$. The constructions that finish on one machine are
the symmetry-prescribed ones, whose search trees are pruned by the prescribed group; the
cost of an orbit-matrix lift grows roughly like $2^{(q-1)\binom{o}{2}}$ in the number of
orbits $o$ and lift modulus $q$, and a commodity machine exhausts only a few tens of bits
of this before wall-clock defeats it. Open-existence graphs, if they exist, are
asymmetric or nearly so, which removes precisely the pruning the tractable regime relies
on. The phenomenon we exploit in this paper is the bright side of the same coin: where a
\emph{seed} of modest symmetry already exists, switching walks into the asymmetric bulk
cheaply; where no graph exists to seed from, there is nothing to walk from. The two
regimes do not overlap, and our results live entirely in the first.

\section{Discussion and limitations}\label{sec:discussion}

\paragraph{What is and is not claimed.} We claim a large, verified family of strongly
regular graphs that are new in the sense of Remark~\ref{rem:caveat}: excluded, by their
automorphism orders, from every identified published enumeration at their parameters. We do
\emph{not} claim a classification, nor that these graphs were never previously written
down by some unpublished switching experiment; at an unclassified parameter neither claim
is available to anyone.

\paragraph{Cell size, and a switching-closed contrast ($\SRG{37,18,8,9}$).} We used $m=4$
throughout, natural because the smallest eigenvalue is $-4$; larger even cells reach further
classes, and a size-$6$ probe already escaped the size-$4$ cascade (Section~\ref{sec:download}).
A sharp contrast is the $\SRG{37,18,8,9}$ catalogue --- the parameter of the long-standing
``$6803$rd graph'' question. Run over the \emph{full} $6802$-graph catalogue, size-$4$ and
size-$8$ GM switching find \emph{no valid cells at all}, while size-$6$ produces $838$ distinct
switches, \emph{every one of which already lies in the catalogue}. The catalogue is therefore
switching-closed, so GM switching cannot exhibit a $6803$rd graph (such a graph, if it exists,
would require the harder regular-two-graph route on $38$ vertices). Notably this closure is
\emph{not} a symmetry effect: the $\SRG{37}$ catalogue is in fact $100\%$ asymmetric
($|\Aut|=1$). The decisive difference from our five parameters is that there the \emph{seed}
set --- a de~Caen lift, or a single SageMath graph --- was \emph{not} switching-closed, so
switching cascaded outward into uncatalogued territory; the $\SRG{37}$ catalogue is already a
complete switching class.

\paragraph{Scope of $|\Aut|$.} The bitset engine is limited to $n<64$. The method extends
verbatim to any unclassified parameter for which low-symmetry seeds are available; the
binding constraint is obtaining such seeds, which for these two parameters came from the
de~Caen lifts.

\paragraph{Future work.} Natural next steps are (i) driving the $\SRG{49,24,11,12}$
cascade to saturation and recording the full automorphism census; (ii) larger cells and
multi-cell GM switching --- a size-$6$ probe already produced $33$ graphs lying outside the
\emph{entire} size-$4$ cascade set of $174{,}209$, proving the reachable region is strictly
larger than the size-$4$ cascade enumerates; and (iii) seeding further unclassified parameters
from fresh low-symmetry constructions --- though this last is itself bounded by the lift wall.
An orbit-matrix (de~Caen-style) attempt at $\SRG{49,18,7,6}$ is instructive: a prescribed
order-$7$ automorphism gave four valid orbit matrices in seconds, yet \emph{none} of them lifted
to a full graph within $1.6\times 10^{7}$ backtracking nodes, so the parameter stayed unopened.
Where a sufficiently asymmetric seed cannot be cheaply obtained --- by a lucky library
construction as for $\SRG{57,24,11,9}$, or a prior lift as for the first three parameters --- the
construction stalls at exactly the wall that stops the open-existence searches
(Section~\ref{sec:ironlaw}). The method therefore opens unclassified parameters readily, but only
where one asymmetric graph is already within reach.

\section{Conclusion}

At realised-but-unclassified parameters, cascading single-cell Godsil--McKay switching
from low-symmetry seeds is a cheap and reliable way to construct very large numbers of new
strongly regular graphs. We produced more than $2.6$ million verified graphs across five
parameters --- $\SRG{50,21,8,9}$, its complement $\SRG{50,28,15,16}$,
$\SRG{49,24,11,12}$, $\SRG{57,24,11,9}$, and its complement $\SRG{57,32,16,20}$ --- the
overwhelming majority asymmetric and
therefore provably outside every symmetry-prescribed census. The governing phenomenon is the
seed-symmetry dichotomy: low-symmetry seeds cascade into the asymmetric bulk, high-symmetry
seeds saturate --- and a single seed of order $|\Aut|=3$ was enough to open the fourth
parameter, while one of order $10$ was not. The same dichotomy clarifies why the open-existence parameters, whose graphs
would be asymmetric, lie beyond the reach of the symmetry-driven constructions that
succeed everywhere else.

\appendix
\section{Reproducibility}\label{app:repro}

The construction is reproduced by: (1) selecting a Brouwer-``$+$'' parameter with $n<64$
and setting the parameters in the switching engine; (2) obtaining low-$|\Aut|$ seeds; (3)
running the cascade of Section~\ref{sec:cascade}, feeding each round's new canonical forms
back as seeds; and (4) verifying every graph for strong regularity and computing $|\Aut|$,
retaining those with $|\Aut|\nmid 6$ (and noting the $|\Aut|=1$ majority). The seed sets,
the switching engine, the cascade driver, the verification scripts, and the cross-checked
parameter database accompany this note in the \texttt{orbit-gen} repository.

\section*{Data availability}
The constructed graphs (in \texttt{graph6} format), the interactive cross-checked database,
this paper, and the construction/verification code are openly available at
\url{https://github.com/tonykoval/srg-database} (static database, draft, and catalogues) and
\url{https://github.com/tonykoval/orbit-gen} (full pipeline). The complete catalogues for
$\SRG{57,24,11,9}$ and its complement ($1{,}122{,}556$ graphs each, $>70$\,MB) are available
on request or regenerable from the seeds and scripts; the full $\SRG{50,21,8,9}$/$\SRG{50,28,15,16}$
catalogues and samples of the larger families are bundled directly. Every graph can be
re-verified by SageMath's \texttt{is\_strongly\_regular()} or nauty.

\begin{thebibliography}{9}
\bibitem{brouwer-table} A.~E.~Brouwer, \emph{Parameters of strongly regular graphs},
  \url{https://aeb.win.tue.nl/graphs/srg/}, tables \texttt{srgtab1-50} and
  \texttt{srgtab51-100}, retrieved 30 May 2026.
\bibitem{bcn} A.~E.~Brouwer, A.~M.~Cohen, A.~Neumaier, \emph{Distance-Regular Graphs},
  Springer, 1989.
\bibitem{bvm-srg} A.~E.~Brouwer, H.~Van~Maldeghem, \emph{Strongly Regular Graphs},
  Encyclopedia of Mathematics and its Applications, Cambridge Univ.\ Press, 2022.
\bibitem{gm1982} C.~D.~Godsil, B.~D.~McKay, \emph{Constructing cospectral graphs},
  Aequationes Math.\ \textbf{25} (1982), 257--268.
\bibitem{ihringer2020} F.~Ihringer, \emph{Switching for small strongly regular graphs and
  related structures}, preprint, arXiv:2012.08390 (2020).
\bibitem{ihringer-switching} F.~Ihringer, A.~Munemasa, \emph{New strongly regular graphs
  from finite geometries via switching}, Linear Algebra Appl.\ \textbf{580} (2019),
  464--474.
\bibitem{behbahani-lam} M.~Behbahani, C.~Lam, \emph{Strongly regular graphs with
  non-trivial automorphisms}, Discrete Math.\ \textbf{311} (2011), 132--144.
\bibitem{decaen} D.~de~Caen, \emph{A lower bound on the probability that two graphs are
  isomorphic, and related lift constructions}; de~Caen-type $\mathbb{Z}_3$ lifts as used
  to generate the seed graphs (\texttt{orbit-gen} project, 2013).
\bibitem{mckay-piperno} B.~D.~McKay, A.~Piperno, \emph{Practical graph isomorphism, II},
  J.\ Symbolic Comput.\ \textbf{60} (2014), 94--112. (nauty/Traces: \texttt{labelg},
  \texttt{dreadnaut}, \texttt{complg}.)
\bibitem{cohen-pasechnik} N.~Cohen, D.~V.~Pasechnik, \emph{Implementing Brouwer's database
  of strongly regular graphs}, Des.\ Codes Cryptogr.\ \textbf{84} (2017), 223--235.
  (SageMath \texttt{sage.graphs.strongly\_regular\_db}.)
\bibitem{sagemath} The Sage Developers, \emph{SageMath, the Sage Mathematics Software
  System} (Version 9.5), 2022, \url{https://www.sagemath.org}.
\bibitem{spence} E.~Spence, \emph{Strongly regular graphs} (online collections), and
  E.~Spence, \emph{Regular two-graphs on 36 vertices}, Linear Algebra Appl.\ \textbf{226--228}
  (1995), 459--497.
\end{thebibliography}

\end{document}
